\centerline{\bf \huge 9 класс}
\centerline{\normalsize{\it Работа рассчитана на \bf{180} минут}}

\problem{1}
Эвена сложила десять подряд идущих натуральных чисел, затем разделила полученную сумму на сумму следующих десяти натуральных чисел. Могло ли у неё получится число $0,8 ?$
%Москва 2017


\problem{2}
Известно, что сумма квадратов корней квадратного уравнения $x^2+b x+ c=0$ равна 2023. Найдите сумму квадратов корней $x^2+b x+c=\frac{1}{2}$.
%Ростов 2023


\problem{3}
В таблицу $4 \times 4$ записали натуральные числа. Могло ли оказаться так, что сумма чисел в каждой следующей строке на 2 больше, чем в предыдущей, а сумма чисел в каждом следующем столбце на 3 больше, чем в предыдущем?
%Ставрополь 2017

\problem{4}
В одну из вершин шестиугольника Скрудж МакДак положил золотую монетку, а в остальных вершинах ничего нет. Каждый день он убирает с одной из вершин шестиугольника произвольное количество монеток и тут же кладёт на соседнюю вершину в шесть раз больше монеток. Если в какойто день Скрудж МакДак сможет добиться того, что во всех вершинах шестиугольника одинаковое число монеток, то ему будет присвоено звание «Великий Селезень». Сможет ли Скрудж МакДак получить это звание?
%Красноярск 2017

\problem{5}
В прямоугольном треугольнике $A B C\left(\angle C=90^{\circ}\right)$ внутри отрезка $B C$ расположена точка $D, C P$ и $C Q$ - высоты треугольников $A B C$ и $A D C$ соответственно. Докажите, что треугольники $A P Q$ и $A D B$ подобны.
%Краснодар 2017